\subsection{A worked dependent connection}
\label{app:worked}

Let the roles be $A,B,C$ and use the script in \eqref{dep:eq:1}, with
\[
 R_A(n,v,j)=\Bool,\qquad R_B(n,v,j)=\mathsf{Fin}(n),
 \qquad R_C(n,v,j)=\mathsf{Vec}(\Bool,n).
\]
Its local types are
\begin{align*}
 \Local{A}{G}
   &=\sum_{n:\Nat_{>0}}\prod_{v:\mathsf{Vec}(\Bool,n)}
                    \prod_{j:\mathsf{Fin}(n)}\Bool,\\
 \Local{B}{G}
   &=\prod_{n:\Nat_{>0}}\sum_{v:\mathsf{Vec}(\Bool,n)}
                    \prod_{j:\mathsf{Fin}(n)}\mathsf{Fin}(n),\\
 \Local{C}{G}
   &=\prod_{n:\Nat_{>0}}\prod_{v:\mathsf{Vec}(\Bool,n)}
                    \sum_{j:\mathsf{Fin}(n)}\mathsf{Vec}(\Bool,n).
\end{align*}
Choose implementations
\begin{align*}
 p_A&=(2,\lambda v.\lambda j.\mathsf{index}(v,j)),\\
 p_B&=\lambda n.(\mathsf{replicate}(n,\mathsf{true}),\lambda j.j),\\
 p_C&=\lambda n.\lambda v.(0_n,v),
\end{align*}
where $0_n:\mathsf{Fin}(n)$ exists because $n$ is positive. The first
interaction selects $2$ and instantiates both receivers at that number.
The next selects Bob's vector $(\mathsf{true},\mathsf{true})$. The third
selects Carol's $0_2$. The final return record is
\[
 \bigl(\mathsf{true},0_2,
       (\mathsf{true},\mathsf{true})\bigr),
\]
indexed by the common trace
$(2,(\mathsf{true},\mathsf{true}),0_2)$.

Neither receiver was required to guess a dependent type from a value it had
not received. The first substitution determines the vector and index types;
the next two substitutions determine the return families. The example
exercises three different owners and all the clauses of the interaction rule.
